\documentclass[11pt,letterpaper]{article}
\usepackage[utf8]{inputenc}
\usepackage[T1]{fontenc}


\usepackage[left=2.4cm,right=2.4cm,top=1.9cm,bottom=1.9cm]{geometry}
\usepackage{array}
\usepackage{booktabs}
\usepackage{xcolor}
\usepackage{fancyhdr}
\usepackage{enumitem}

\usepackage[hidelinks]{hyperref}

\definecolor{gris}{gray}{0.35}
\pagestyle{fancy}
\fancyhf{}
\fancyhead[L]{\footnotesize\color{gris} Acta de aprobación --- SRS v1.0.0 BIOPET}
\fancyhead[R]{\footnotesize\color{gris} ISO/IEC/IEEE 29148:2018}
\fancyfoot[C]{\footnotesize\color{gris} \thepage}
\renewcommand{\headrulewidth}{0.4pt}

\setlist[itemize]{leftmargin=1.3em,itemsep=1pt,topsep=2pt}

\begin{document}

\begin{center}
	{\footnotesize UNIVERSIDAD TÉCNICA ESTATAL DE QUEVEDO\\[1pt]
	FACULTAD DE CIENCIAS DE LA COMPUTACIÓN --- CARRERA DE SOFTWARE\\[1pt]
	Proyecto de Fin de Carrera --- Aplicaciones Web --- Periodo académico 2025--2026}\\[10pt]
	\rule{\textwidth}{0.8pt}\\[8pt]
	{\large\bfseries ACTA DE APROBACIÓN DEL DOCENTE-DIRECTOR}\\[3pt]
	{\normalsize Especificación de Requisitos de Software (SRS)}\\[6pt]
	\rule{\textwidth}{0.8pt}
\end{center}

\vspace{0pt}

\section*{1. Identificación del documento aprobado}

\renewcommand{\arraystretch}{1.08}
\noindent\begin{tabular}{@{}>{\bfseries}p{4.6cm}p{10.4cm}@{}}
	\toprule
	Proyecto & BIOPET --- Sistema web de gestión veterinaria \\
	Documento & Especificación de Requisitos de Software (SRS) \\
	Versión & v1.0.0 --- Entrega Final \\
	Equipo & BMT \\
	Integrantes & Fred Adrián Beltrán Montiel (ORCID 0009-0007-8303-2137) \newline Jaime Josué Mariscal Cabrera (ORCID 0009-0007-2206-3941) \newline Zaida Melissa Taipe Mora (ORCID 0009-0000-1227-3258) \\
	Repositorio & \url{https://github.com/Grinjoww/ENTREGA-FINAL-BIOPET} \\
	Corpus & 44 requisitos (25 funcionales y 19 no funcionales) \\
	Fecha de aprobación & 12 de septiembre de 2026 \\
	\bottomrule
\end{tabular}

\section*{2. Declaración de aprobación}

Apruebo la Especificación de Requisitos de Software (SRS v1.0.0) del proyecto BIOPET como especificación válida de la Entrega Final del Proyecto de Fin de Carrera de la asignatura Aplicaciones Web.

Dejo constancia expresa de que esta aprobación recae sobre el documento en su calidad de especificación, y no constituye una certificación de que el sistema especificado carezca de defectos.

\section*{3. Verificación realizada}

Contrasté el documento contra el repositorio declarado en su portada y constato lo siguiente.

\begin{itemize}
	\item Los 44 identificadores de requisito del documento coinciden exactamente con los 44 registrados en \texttt{docs/trazabilidad/matriz.csv}, sin identificadores huérfanos en ninguna de las dos direcciones.
	\item Los 40 bloques de requisito aplican la plantilla completa, con criterio de aceptación observable, método de verificación, trazabilidad, prioridad, \emph{rationale} y estado. En la revisión anterior únicamente tres requisitos declaraban criterio de aceptación.
	\item El vocabulario de estados quedó cerrado a tres valores y se aplica sin excepción en todo el documento.
	\item Las secciones de interfaces externas, matriz de permisos por rol y operación, y estados del dominio, ausentes en la revisión anterior, están construidas a partir del código del propio repositorio.
	\item Las historias de usuario, los casos de uso y las clases de prueba citadas como método de verificación existen en el repositorio.
\end{itemize}

\vspace{-4pt}
\section*{4. Observación expresa sobre el estado del sistema}

Apruebo con la siguiente observación, que el equipo debe atender y que afecta al estado del sistema, no a la validez de este documento.

El sistema presenta un defecto de aislamiento de datos que el propio documento declara con precisión en REQ-F-006, REQ-F-013, REQ-NF-015, la sección 2.8.5 y la sección 7.3, y que verifiqué vigente en la rama principal del repositorio.

El método de listado de consultas médicas evalúa el rol del solicitante, pero ambas ramas de la condición devuelven la misma consulta sin filtro por propietario, y el repositorio correspondiente no declara el método de filtrado que sí existe en los repositorios de citas y de vacunas.

En consecuencia, un usuario autenticado con rol de dueño de mascota puede listar consultas médicas de mascotas ajenas mediante \texttt{GET /api/consultas}.

Ese comportamiento incumple el criterio 3 de REQ-F-006, requisito de prioridad \emph{Must}, por lo que considero correcto que el equipo lo haya declarado en estado \emph{implementado} y no \emph{verificado}.

La corrección requiere añadir al repositorio de consultas un método de filtrado por propietario y emplearlo en la rama correspondiente del servicio, replicando el patrón del módulo de mascotas.

\vspace{-4pt}
\section*{5. Recomendaciones para trabajo futuro}

\begin{itemize}
	\item El sistema no dispone de mecanismo de recuperación de contraseña, pese a que el registro de usuarios es público. Un usuario que olvide su credencial queda sin vía de acceso especificada.
	\item El comportamiento ante la indisponibilidad del servicio de caché no está controlado: la consulta a la lista de revocación de tokens no captura la excepción de conexión, por lo que una caída de ese servicio deriva en error \texttt{500} en toda petición autenticada, sin que exista una decisión declarada de rechazar o aceptar en ese escenario.
\end{itemize}

\vspace{-2pt}
\section*{6. Mérito registrado}

Valoro expresamente la honestidad metodológica del documento: el equipo declara el defecto de aislamiento en cinco lugares distintos del texto, rebaja por ello el estado de un requisito \emph{Must} que en la versión anterior figuraba como verificado, y enumera sus limitaciones abiertas sin cerrarlas con datos no comprobables. Registro esa conducta como mérito del equipo.

\vspace{6pt}

\noindent\rule{\textwidth}{0.4pt}

\vspace{4pt}

\noindent\textbf{7. Suscripción}

\vspace{18pt}

\noindent\begin{tabular}{@{}p{10.5cm}@{}}
	\hrulefill \\[3pt]
	\textbf{Ing. Gleiston Cicerón Guerrero Ulloa, Ph.D.} \\
	Docente-director del Proyecto de Fin de Carrera \\
	Asignatura: Aplicaciones Web \\
	Universidad Técnica Estatal de Quevedo \\
\end{tabular}

\vspace{8pt}

\noindent{\footnotesize\color{gris} Documento suscrito mediante firma electrónica; su validez se verifica con el certificado digital incorporado a este archivo.}

\end{document}
